\nonstopmode{}
\documentclass[a4paper]{book}
\usepackage[times,inconsolata,hyper]{Rd}
\usepackage{makeidx}
\makeatletter\@ifl@t@r\fmtversion{2018/04/01}{}{\usepackage[utf8]{inputenc}}\makeatother
% \usepackage{graphicx} % @USE GRAPHICX@
\makeindex{}
\begin{document}
\chapter*{}
\begin{center}
{\textbf{\huge Package `coorsim'}}
\par\bigskip{\large \today}
\end{center}
\ifthenelse{\boolean{Rd@use@hyper}}{\hypersetup{pdftitle = {coorsim: Detects Coordinated Posting Behavior Based on Co-Similarity}}}{}
\ifthenelse{\boolean{Rd@use@hyper}}{\hypersetup{pdfauthor = {Thiele Daniel}}}{}
\begin{description}
\raggedright{}
\item[Type]\AsIs{Package}
\item[Title]\AsIs{Detects Coordinated Posting Behavior Based on Co-Similarity}
\item[Version]\AsIs{0.0.7}
\item[Date]\AsIs{2026-06-28}
\item[Description]\AsIs{R package for detection of coordinated posting behavior and coordinated social media manipulation. 
Detects groups of social media accounts that post semantically similar content within short time frames, by comparing cosine similarities of post embeddings. 
Provides a comprehensive set of functions to generate embeddings (via reticulate, pytorch), to efficiently compute co-similarities (via C++ and tbb),
to detect coordinated communites, automatically label communities with locally hosted LLMs (via ollama), and to plot results.}
\item[License]\AsIs{GPL (>= 3)}
\item[Encoding]\AsIs{UTF-8}
\item[Imports]\AsIs{assertthat, cli, data.table, dplyr, igraph, lubridate, Matrix,
Rcpp (>= 1.0.12), reticulate, stats, stringr, hdf5r, furrr,
proxyC, text2map, pbapply, jsonlite, purrr, tibble}
\item[LinkingTo]\AsIs{Rcpp, RcppArmadillo (>= 0.7.600.1.0)}
\item[RoxygenNote]\AsIs{7.3.2}
\item[Depends]\AsIs{R (>= 4.3.0)}
\item[NeedsCompilation]\AsIs{yes}
\item[SystemRequirements]\AsIs{GNU make}
\item[Suggests]\AsIs{emoji, future, httr, ggplot2, ggraph, ggforce, ggrepel,
ggtext, parallel, patchwork, quanteda, rollama, stringdist,
stringi, testthat (>= 3.0.0), tidygraph, viridis, xgboost,
mltools, ineq}
\item[Config/testthat/edition]\AsIs{3}
\item[LazyData]\AsIs{true}
\item[Author]\AsIs{Thiele Daniel [cre, aut]}
\item[Maintainer]\AsIs{Thiele Daniel }\email{daniel.thiele@fu-berlin.de}\AsIs{}
\item[Archs]\AsIs{x64}
\end{description}
\Rdcontents{Contents}
\HeaderA{.reconstruct\_description\_v2}{Reconstruct a 'description' value from the flat desc\_* fields (v0.5 schema)}{.reconstruct.Rul.description.Rul.v2}
%
\begin{Description}
Reads desc\_lang\_topics / desc\_named\_entities / desc\_patterns /
desc\_tone\_style / desc\_striking from a parsed JSON list and returns either
a JSON string (keep\_json = TRUE) or a space-joined sentence string
(keep\_json = FALSE), matching the output contract of .process\_comm\_description().
\end{Description}
%
\begin{Usage}
\begin{verbatim}
.reconstruct_description_v2(parsed_json, keep_json = FALSE)
\end{verbatim}
\end{Usage}
%
\begin{Arguments}
\begin{ldescription}
\item[\code{parsed\_json}] list. A parsed JSON response from the v0.5 schema.

\item[\code{keep\_json}] logical. If TRUE returns a JSON object string; else plain text.
\end{ldescription}
\end{Arguments}
%
\begin{Details}
The mapping to the canonical sentence numbering is:
sentence\_1 <- desc\_lang\_topics
sentence\_2 <- desc\_named\_entities
sentence\_3 <- desc\_patterns
sentence\_4 <- desc\_tone\_style
sentence\_5 <- desc\_striking
\end{Details}
%
\begin{Value}
character(1)
\end{Value}
\HeaderA{augment\_groups\_data}{Augment Group Data with Post and User Information}{augment.Rul.groups.Rul.data}
%
\begin{Description}
This function merges additional post and user data into a grouped data structure, `groups\_data`.
It supports renaming columns for consistency, de-duplication, and optional community annotation in the similarity table.
\end{Description}
%
\begin{Usage}
\begin{verbatim}
augment_groups_data(
  groups_data,
  post_data,
  user_data,
  post_id = NULL,
  account_id = NULL,
  other_post_vars = NULL,
  other_user_vars = NULL,
  sim_dt_community = TRUE,
  verbose = FALSE
)
\end{verbatim}
\end{Usage}
%
\begin{Arguments}
\begin{ldescription}
\item[\code{groups\_data}] A list with at least one element `node\_list` (a `data.table` containing node information).
Optionally, it may also contain `post\_data` and `sim\_dt`, which will be augmented.

\item[\code{post\_data}] A `data.table` containing additional post data to be merged.

\item[\code{user\_data}] A `data.table` containing user data to be merged.

\item[\code{post\_id}] Optional. Column name in `post\_data` to be renamed to `'post\_id'` for joining.

\item[\code{account\_id}] Optional. Column name in `user\_data` and `groups\_data` to be renamed to `'account\_id'` for joining.

\item[\code{other\_post\_vars}] Optional. Additional columns to retain from `post\_data`.

\item[\code{other\_user\_vars}] Optional. Additional columns to retain from `user\_data`.

\item[\code{sim\_dt\_community}] Logical. If `TRUE`, community identifiers from `node\_list` will be joined onto `sim\_dt`,
matching on `account\_id` and `account\_id\_y`. The resulting `sim\_dt` will be reordered
to follow the structure: x-side → similarity → y-side → parameters.

\item[\code{verbose}] Logical. If `TRUE`, provides detailed output about each processing step.
\end{ldescription}
\end{Arguments}
%
\begin{Details}
This function standardizes column names, selects relevant fields, and performs efficient `data.table` joins.
User-level data is prefixed with \code{"account\_"} before merging. If \code{sim\_dt\_community = TRUE}, community IDs
are merged into `sim\_dt` and the column order is standardized to place the x-side account and post information first,
followed by similarity metrics, y-side information, and parameter columns.
\end{Details}
%
\begin{Value}
A list, `groups\_data`, with the following possible modifications:
\begin{itemize}

\item{} \code{post\_data}: augmented with user information and additional post columns, if present.
\item{} \code{node\_list}: augmented with user data and optionally sampled content from post data.
\item{} \code{sim\_dt}: (if \code{sim\_dt\_community = TRUE}) augmented with \code{community} and \code{community\_y}
for \code{account\_id} and \code{account\_id\_y}, and columns reordered accordingly.

\end{itemize}

\end{Value}
\HeaderA{check\_backend}{Check coorsim Backend Configuration}{check.Rul.backend}
%
\begin{Description}
Displays diagnostic information about the coorsim Python backend setup.
\end{Description}
%
\begin{Usage}
\begin{verbatim}
check_backend()
\end{verbatim}
\end{Usage}
%
\begin{Value}
Invisibly returns a list with backend configuration details.
\end{Value}
%
\begin{Examples}
\begin{ExampleCode}
## Not run: 
  coorsim::check_backend()

## End(Not run)
\end{ExampleCode}
\end{Examples}
\HeaderA{check\_gpu}{CUDA-detection only, just for onnxruntime-gpu}{check.Rul.gpu}
%
\begin{Description}
CUDA-detection only, just for onnxruntime-gpu
\end{Description}
%
\begin{Usage}
\begin{verbatim}
check_gpu()
\end{verbatim}
\end{Usage}
\HeaderA{clean\_json\_strings}{Clean and repair JSON-like strings generated by LLMs}{clean.Rul.json.Rul.strings}
%
\begin{Description}
This function attempts to sanitize JSON-like character strings returned by large language models (LLMs),
preparing them for parsing via `jsonlite::fromJSON()`. It extracts the first substring that resembles
a JSON object, replaces single quotes with double quotes (when safe), escapes backslashes,
and ensures the string is properly enclosed in curly braces.
\end{Description}
%
\begin{Usage}
\begin{verbatim}
clean_json_strings(json_strings)
\end{verbatim}
\end{Usage}
%
\begin{Arguments}
\begin{ldescription}
\item[\code{json\_strings}] A character vector of JSON-like strings (often malformed) generated by an LLM.
\end{ldescription}
\end{Arguments}
%
\begin{Details}
This function is not guaranteed to fix all malformed JSON but is useful for cleaning common
LLM formatting issues. It does not validate the final JSON structure—only prepares it for parsing.
\end{Details}
%
\begin{Value}
A character vector of cleaned strings, formatted to increase the likelihood of successful parsing
with `jsonlite::fromJSON()`.
\end{Value}
%
\begin{SeeAlso}
[jsonlite::fromJSON()]
\end{SeeAlso}
%
\begin{Examples}
\begin{ExampleCode}
# Example 1: Fix single quotes
raw_json <- "{'label':'Test','description':'A test case','lang':'en'}"
clean_json_strings(raw_json)

# Example 2: Handle backslashes in paths
messy_json <- "{'label':'Path','description':'Folder path is C:\\\\Users\\\\','lang':'en'}"
clean_json_strings(messy_json)

# Example 3: Add missing braces
no_braces <- "'label':'Oops','lang':'en'"
clean_json_strings(no_braces)

\end{ExampleCode}
\end{Examples}
\HeaderA{coorsim\_detect\_groups}{Detect Communities of Coordinated Accounts}{coorsim.Rul.detect.Rul.groups}
%
\begin{Description}
This function identifies communities of coordinated accounts in a co-similarity network. 
It supports two community detection algorithms—the classic "Louvain" algorithm, the "Leiden" algorithm, the "label\_propagation" algorithm,
and the Focal Structures Algorithm "FSA\_V" proposed by Weber \& Neumann (2021), based on Sen et al. (2016). 
It can filter communities based on a minimum edge weight threshold and filter the network, and all other output by minimum community size.
\end{Description}
%
\begin{Usage}
\begin{verbatim}
coorsim_detect_groups(
  simdt,
  user_data,
  account_id = "account_id",
  account_name = "account_name",
  other_user_vars = NULL,
  edge_weight = NULL,
  cluster_method = "label_prop",
  resolution = 1,
  theta = 0.3,
  return_post_dt = TRUE,
  verbose = TRUE,
  seed = 42,
  min_comm_size = NULL,
  ...
)
\end{verbatim}
\end{Usage}
%
\begin{Arguments}
\begin{ldescription}
\item[\code{simdt}] A `data.table` containing the network edges, with columns for account IDs 
(`account\_id`, `account\_id\_y`) and optional columns for `post\_id` and `time`.

\item[\code{user\_data}] A `data.table` of account-level data with at least one column for `account\_id`. 
Additional columns can be specified to be included in the output.

\item[\code{account\_id}] Name of the column in `user\_data` that identifies accounts. Defaults to "account\_id".

\item[\code{account\_name}] Name of the column in `user\_data` for account names or labels. Defaults to "account\_name".

\item[\code{other\_user\_vars}] Optional character vector of additional column names in `user\_data` 
to include in the output.

\item[\code{edge\_weight}] Optional numeric threshold for filtering the edge list. 
Only edges with a weight greater than or equal to this value will be included.

\item[\code{cluster\_method}] Character, the algorithm to use for community detection. 
Options are "louvain", "leiden", "label\_prop" or "FSA\_V". Defaults to "label\_prop".

\item[\code{resolution}] Numeric resolution parameter for the Louvain algorithm. 
Controls the size of detected communities. Higher values create smaller communities. Defaults to `1`.

\item[\code{theta}] Numeric, specifies the threshold for the FSA\_V algorithm to control 
the inclusion of edges. Higher values prioritize stronger edges in community construction. Defaults to `0.7`.

\item[\code{return\_post\_dt}] Logical, whether to return a `data.table` with post data merged with community labels. 
Defaults to `TRUE`.

\item[\code{verbose}] Logical, whether to display progress messages. Defaults to `TRUE`.

\item[\code{seed}] Numeric, the random seed to ensure reproducibility. Defaults to `42`.

\item[\code{min\_comm\_size}] Integer. Minimum community size for filtering. Default is `NULL`.

\item[\code{...}] Additional parameters passed to the clustering functions.
\end{ldescription}
\end{Arguments}
%
\begin{Value}
A named list containing:
\begin{description}

\item[\code{graph}] An `igraph` object representing the network graph created from the filtered edge list.
\item[\code{communities}] An `igraph::communities` object representing the detected communities.
\item[\code{edge\_list}] A filtered `data.table` of edges used for community detection.
\item[\code{node\_list}] A `data.table` with node membership information, including account ID, community assignment, and algorithm used.
\item[\code{post\_data}] A `data.table` containing post-level data with community information, only returned if \code{return\_post\_dt} is `TRUE`.

\end{description}

\end{Value}
\HeaderA{coorsim\_match\_overlaps}{Match Overlapping Posts by Time Window}{coorsim.Rul.match.Rul.overlaps}
%
\begin{Description}
This function matches overlapping posts within a specified time window, removes duplicated pairs,
and groups overlaps by `post\_id` using data.table functions.
\end{Description}
%
\begin{Usage}
\begin{verbatim}
coorsim_match_overlaps(data)
\end{verbatim}
\end{Usage}
%
\begin{Arguments}
\begin{ldescription}
\item[\code{data}] A data.table containing the columns `post\_id`, `account\_id`, `time`, `end\_time`, and `time\_dup`.
\end{ldescription}
\end{Arguments}
%
\begin{Value}
A named list where each name corresponds to a `post\_id` and each element is a vector of post IDs that overlap within the specified time window.
\end{Value}
\HeaderA{coorsim\_prepare\_data}{Prepare Data for Co-Similarity Analysis}{coorsim.Rul.prepare.Rul.data}
%
\begin{Description}
This function processes and prepares social media post data for correlation similarity analysis.
It ensures that required columns are present, standardizes time formats, and optionally integrates
embeddings either from an in-memory matrix or an `.h5` file.
\end{Description}
%
\begin{Usage}
\begin{verbatim}
coorsim_prepare_data(
  data,
  embeddings = NULL,
  time_window = 60,
  post_id = NULL,
  account_id = NULL,
  time = NULL,
  content = NULL,
  verbose = TRUE
)
\end{verbatim}
\end{Usage}
%
\begin{Arguments}
\begin{ldescription}
\item[\code{data}] A `data.frame` or `data.table` containing post data. The table must include at least
`post\_id`, `account\_id`, `time`, and `content`, either directly or via specified alternative column names.

\item[\code{embeddings}] Embeddings for the posts, which can be provided in **three formats**:
\begin{itemize}

\item{} **A matrix**: Row names must correspond to `post\_id` values.
\item{} **A valid `.h5` file path**: The file must contain a `metadata/post\_id` dataset.
\item{} **NULL (default)**: No embeddings are used.

\end{itemize}


\item[\code{time\_window}] Integer. The time window (in seconds) within which similar posts are detected. Default is `60`.

\item[\code{post\_id}] Character. Alternative column name for `post\_id`, if not present in `data`.

\item[\code{account\_id}] Character. Alternative column name for `account\_id`, if not present in `data`.

\item[\code{time}] Character. Alternative column name for `time`, if not present in `data`.

\item[\code{content}] Character. Alternative column name for `content`, if not present in `data`.

\item[\code{verbose}] Logical. If `TRUE`, displays progress messages. Default is `TRUE`.
\end{ldescription}
\end{Arguments}
%
\begin{Details}
This function performs the following preprocessing steps:
\begin{itemize}

\item{} Ensures that `post\_id`, `account\_id`, `time`, and `content` exist in `data`,
renaming alternative column names if provided.
\item{} Drops duplicate posts based on `post\_id`.
\item{} Standardizes `time` to **Unix timestamps (UTC)** for cross-platform compatibility.
\item{} Sorts posts in ascending order by `time`.
\item{} If **embeddings are provided as a matrix**, retains only posts present in `rownames(embeddings)`.
\item{} If **embeddings are provided as an `.h5` file**, verifies that `metadata/post\_id` exists
before extracting relevant post IDs.

\end{itemize}


The function also filters out invalid time formats and ensures efficient memory usage by removing the `content` column
from the final output.
\end{Details}
%
\begin{Value}
A processed `data.table` ready for correlation similarity analysis.
\end{Value}
\HeaderA{detect\_cosimilarity}{Detect Co-Similarities}{detect.Rul.cosimilarity}
%
\begin{Description}
This function identifies highly similar social media posts shared within a specified time window,
aiming to detect coordinated social media manipulation. It uses either text-based sparse document-feature matrices
or numerical embeddings to compute pairwise similarity scores.
\end{Description}
%
\begin{Usage}
\begin{verbatim}
detect_cosimilarity(
  data,
  embeddings = NULL,
  time_window = 60,
  min_simil = 0.9,
  min_participation = 3,
  post_id = NULL,
  account_id = NULL,
  time = NULL,
  content = NULL,
  verbose = TRUE,
  method = "cosine",
  remove_loops = TRUE,
  parallel = TRUE,
  n_threads = NULL,
  subset_emb = FALSE,
  embeddings_name = NULL
)
\end{verbatim}
\end{Usage}
%
\begin{Arguments}
\begin{ldescription}
\item[\code{data}] A `data.frame` or `data.table` containing social media data. Must include columns for post IDs, account IDs, timestamps, and content.

\item[\code{embeddings}] Optional. Embeddings for the posts, which can be provided in the following formats:
\begin{itemize}

\item{} **A matrix**: Row names must correspond to `post\_id` values.
\item{} **A valid `.h5` file path**: The file must contain a `metadata/post\_id` dataset.
\item{} **NULL (default)**: Uses text-based similarity instead of embeddings.

\end{itemize}


\item[\code{time\_window}] Integer. Specifies the time window (in seconds) within which to detect similar posts. Default is `60`.

\item[\code{min\_simil}] Numeric. The minimum similarity threshold to consider posts as similar. Default is `0.9`.

\item[\code{min\_participation}] Integer. The minimum number of participations required to consider an account as part of coordinated activity. Default is `3`.

\item[\code{post\_id}] Optional. Character string specifying the column name for post IDs if different from "post\_id".

\item[\code{account\_id}] Optional. Character string specifying the column name for account IDs if different from "account\_id".

\item[\code{time}] Optional. Character string specifying the column name for timestamps if different from "time".

\item[\code{content}] Optional. Character string specifying the column name for post content if different from "content".

\item[\code{verbose}] Logical. If `TRUE`, displays progress messages. Default is `TRUE`.

\item[\code{method}] Character. Specifies the similarity method to use (e.g., "cosine"). Default is "cosine".

\item[\code{remove\_loops}] Logical. If `TRUE`, removes self-similar posts from the same account. Default is `TRUE`.

\item[\code{parallel}] Logical. If `TRUE`, enables parallel computation for similarity calculations. Default is `TRUE`.

\item[\code{n\_threads}] Optional. Integer specifying the number of threads for parallel computation. If `NULL`, defaults to available cores minus one.

\item[\code{subset\_emb}] Logical. Should the embedding matrix be subsetted to observations in data? Slower but more memory-friendly. Default is `FALSE`.

\item[\code{embeddings\_name}] Character. The name of the embedding model used to be stored in 'sim\_dt'. Default is `NULL`.
\end{ldescription}
\end{Arguments}
%
\begin{Details}
The function follows four main steps:
\begin{enumerate}

\item{} **Preprocess Data**: Standardizes column names, ensures valid timestamps, and prepares embeddings if provided.
\item{} **Match Overlapping Posts**: Identifies posts published within the specified `time\_window`.
\item{} **Compute Similarities**: 
\begin{itemize}

\item{} If `embeddings` are `NULL`, constructs a sparse document-feature matrix and calculates similarity scores.
\item{} If `embeddings` are provided, queries and computes pairwise similarities using optimized C++ functions.

\end{itemize}

\item{} **Filter and Aggregate**: Joins account and timestamp data, filters accounts by `min\_participation`, and removes redundant pairs.

\end{enumerate}

\end{Details}
%
\begin{Value}
A `data.table` containing pairs of similar posts with the following columns:
\begin{itemize}

\item{} `post\_id` - The ID of the first post in the pair.
\item{} `post\_id\_y` - The ID of the second post in the pair.
\item{} `similarity` - The computed similarity score.
\item{} `account\_id` - The account that posted `post\_id`.
\item{} `account\_id\_y` - The account that posted `post\_id\_y`.
\item{} `time` - The timestamp of `post\_id`.
\item{} `time\_y` - The timestamp of `post\_id\_y`.
\item{} `time\_diff` - The time difference between the two posts (in seconds).

\end{itemize}

\end{Value}
\HeaderA{filter\_groups\_data}{Filter a detected group structure based on various criteria}{filter.Rul.groups.Rul.data}
%
\begin{Description}
This function filters a detected group structure (`groups\_data`) by a wide range of optional thresholds
and selection criteria, including account-level, post-level, and similarity-based filters. It supports
optional re-detection of groups after filtering similarity pairs (`sim\_dt`), and returns a cleaned, 
consistent `groups\_data` object.
\end{Description}
%
\begin{Usage}
\begin{verbatim}
filter_groups_data(
  groups_data,
  by_col = NULL,
  by_val = NULL,
  by_condition = "equal",
  edge_weight = NULL,
  time_window = NULL,
  min_simil = NULL,
  min_comm_size = NULL,
  min_comp_size = NULL,
  min_degree = NULL,
  min_participation = NULL,
  communities_index = NULL,
  quantile_accountwise = NULL,
  quantile_postwise = NULL,
  top_n_accountwise = NULL,
  top_n_postwise = NULL,
  stringsimil_cutoff = NULL,
  stringsimil_method = "cosine",
  stringsimil_ngram = 3,
  nthread = 10,
  rerun_detect_groups = FALSE,
  verbose = TRUE,
  ...
)
\end{verbatim}
\end{Usage}
%
\begin{Arguments}
\begin{ldescription}
\item[\code{groups\_data}] A list returned by `coorsim\_detect\_groups()` or `augment\_groups\_data()`.

\item[\code{by\_col}] Column name in `post\_data` to filter by (e.g. `"lang"`, `"country"`).

\item[\code{by\_val}] Value(s) to keep in `by\_col`.

\item[\code{by\_condition}] Condition to compare values in `by\_col` to the values in `by\_val`. Options: 'equal', 'greater'. 'smaller'
Default is 'equal'.

\item[\code{edge\_weight}] Minimum edge weight to retain in the similarity network.

\item[\code{time\_window}] Maximum time difference allowed between posts (in seconds).

\item[\code{min\_simil}] Minimum cosine similarity allowed between post embeddings.

\item[\code{min\_comm\_size}] Minimum number of accounts in a community.

\item[\code{min\_comp\_size}] Minimum number of nodes in a graph component.

\item[\code{min\_degree}] Minimum node degree.

\item[\code{min\_participation}] Minimum number of post-level similarity relations per account.

\item[\code{communities\_index}] Vector of community indices to keep.

\item[\code{quantile\_accountwise}] Drop communities below a given quantile of account size.

\item[\code{quantile\_postwise}] Drop communities below a given quantile of post size.

\item[\code{top\_n\_accountwise}] Keep the top-n largest communities by account count.

\item[\code{top\_n\_postwise}] Keep the top-n largest communities by post count.

\item[\code{stringsimil\_cutoff}] Optional. Additional filter: minimum geometric mean of cosine similarity and string similarity.

\item[\code{stringsimil\_method}] String similarity method passed to `stringdist::stringsim()` (e.g. `"cosine"`, `"jaccard"`).

\item[\code{stringsimil\_ngram}] Integer n-gram size used in string similarity.

\item[\code{nthread}] Number of threads used in `stringsim` filtering.

\item[\code{rerun\_detect\_groups}] Logical. Whether to re-run group detection after filtering similarity pairs.

\item[\code{verbose}] Logical. Whether to print informative messages.

\item[\code{...}] Additional arguments passed to `coorsim\_detect\_groups()` if rerun is triggered.
\end{ldescription}
\end{Arguments}
%
\begin{Value}
A filtered `groups\_data` list with updated `\$filter` slot for reproducibility.
\end{Value}
%
\begin{SeeAlso}
[coorsim\_detect\_groups()], [augment\_groups\_data()]
\end{SeeAlso}
\HeaderA{filter\_ntile}{Filter Data by Quantile Threshold}{filter.Rul.ntile}
%
\begin{Description}
This function filters a data frame or data.table by retaining rows where the specified variable is 
greater than or equal to a given quantile threshold. It supports filtering by the variable's name 
or by its index position.
\end{Description}
%
\begin{Usage}
\begin{verbatim}
filter_ntile(df, var_field, probs)
\end{verbatim}
\end{Usage}
%
\begin{Arguments}
\begin{ldescription}
\item[\code{df}] A data frame or data.table to filter.

\item[\code{var\_field}] Character or integer indicating the column to filter by. If character, it should match the column name; 
if integer, it should match the column index.

\item[\code{probs}] Numeric value between 0 and 1 indicating the quantile threshold for filtering. Rows with values 
in `var\_field` greater than or equal to this quantile are retained.
\end{ldescription}
\end{Arguments}
%
\begin{Value}
The filtered data frame or data.table containing rows where the specified variable is at or above the specified quantile.
\end{Value}
\HeaderA{flatten\_description}{Flatten a JSON Description String to Plain Text}{flatten.Rul.description}
%
\begin{Description}
Parses a JSON string of the form `\{"sentence\_1": "...", "sentence\_2": "...", ...\}`
and concatenates the values into a single space-separated string. Non-JSON strings
and `NA` values are returned as-is.
\end{Description}
%
\begin{Usage}
\begin{verbatim}
flatten_description(x)
\end{verbatim}
\end{Usage}
%
\begin{Arguments}
\begin{ldescription}
\item[\code{x}] Character vector of JSON strings (or plain text) to flatten.
\end{ldescription}
\end{Arguments}
%
\begin{Value}
A character vector of the same length as \code{x}.
\end{Value}
\HeaderA{get\_community\_metrics}{Compute and aggregate multi-domain community-level metrics}{get.Rul.community.Rul.metrics}
%
\begin{Description}
This function computes a comprehensive set of community-level metrics based on the 
input `groups\_data` object, which includes graph, similarity, and account data. It 
extracts and aggregates network structure metrics, similarity and temporal coordination 
statistics, and account-based indicators. Optionally, it also computes lexical metrics 
from post and account-level text content.
\end{Description}
%
\begin{Usage}
\begin{verbatim}
get_community_metrics(
  groups_data,
  account_description = NULL,
  account_creation_date = NULL,
  follower = NULL,
  following = NULL,
  content_stats = TRUE,
  zstd = FALSE,
  verbose = T
)
\end{verbatim}
\end{Usage}
%
\begin{Arguments}
\begin{ldescription}
\item[\code{groups\_data}] A list containing the following elements:
- `graph`: an igraph object representing the network.
- `communities`: an igraph community object (e.g. from `cluster\_louvain()`).
- `node\_list`: a data.table with account-level metadata.
- `sim\_dt`: a data.table with coordinated post pairs and similarity info.
- `post\_data` (optional): a data.table with post-level content and metadata.

\item[\code{account\_description}] Name of the column containing account descriptions (default `NULL`). 
If not present under standard name, the function will attempt to use this.

\item[\code{account\_creation\_date}] Name of the column with account creation timestamps (default `NULL`).

\item[\code{follower}] Name of the follower count column (default `NULL`).

\item[\code{following}] Name of the following count column (default `NULL`).

\item[\code{content\_stats}] Logical; if `TRUE`, compute and aggregate lexical content metrics (default `TRUE`).

\item[\code{zstd}] Logical; if `TRUE` will z-standardize all columns (default `FALSE`).

\item[\code{verbose}] Logical; if `TRUE`, shows progress messages using the `cli` package (default `TRUE`).
\end{ldescription}
\end{Arguments}
%
\begin{Details}
The function internally computes:
- Network statistics from subgraphs induced by communities
- Similarity and time-difference metrics from coordinated post pairs
- Posting temporal entropy and burstiness (via 3-hour bins)
- Optional text metrics from post and account descriptions (if `content\_stats = TRUE`)

It automatically standardizes column names for description, creation date,
follower/following if alternate column names are supplied.
\end{Details}
%
\begin{Value}
A data.table where each row represents one community and columns contain
metrics for:
- Network topology (e.g., degree, centrality, density)
- Similarity structure (e.g., average pairwise similarity within communities)
- Temporal coordination (e.g., entropy of posting over time)
- Account-level metadata (e.g., account age, follower ratio, post count)
- Optionally, post and bio content elaborateness metrics (if `content\_stats = TRUE`)
\end{Value}
%
\begin{SeeAlso}
[coorsim::get\_content\_metrics()], [coorsim::z\_standardize()]
\end{SeeAlso}
\HeaderA{get\_content\_metrics}{Compute and aggregate lexical content metrics per community}{get.Rul.content.Rul.metrics}
%
\begin{Description}
Computes content-level elaborateness metrics using quanteda, and
aggregates them per community as mean and sd. ID column is user-defined.
\end{Description}
%
\begin{Usage}
\begin{verbatim}
get_content_metrics(
  dt,
  content_col = "text",
  id_col = "post_id",
  prefix = "c_"
)
\end{verbatim}
\end{Usage}
%
\begin{Arguments}
\begin{ldescription}
\item[\code{dt}] A data.table with a content column, community column, and an ID column

\item[\code{content\_col}] Name of the text column (default: "text")

\item[\code{id\_col}] Name of the unique ID column (e.g., "post\_id", "account\_id")

\item[\code{prefix}] Prefix to add to all metric columns in output (default: "c\_")
\end{ldescription}
\end{Arguments}
%
\begin{Value}
data.table with one row per community and aggregated metrics
\end{Value}
\HeaderA{get\_embeddings}{Retrieve Text Embeddings Using Hugging Face Model via Python Script}{get.Rul.embeddings}
%
\begin{Description}
Retrieve Text Embeddings Using Hugging Face Model via Python Script
\end{Description}
%
\begin{Usage}
\begin{verbatim}
get_embeddings(
  data,
  post_id = NULL,
  time = NULL,
  content = NULL,
  model_name = "sentence-transformers/paraphrase-multilingual-MiniLM-L12-v2",
  batch_size = 16L,
  max_length = 512L,
  use_fp16 = TRUE,
  trust_remote_code = TRUE,
  verbose = TRUE,
  uv_cache_dir = NULL,
  models_dir = NULL
)
\end{verbatim}
\end{Usage}
%
\begin{Arguments}
\begin{ldescription}
\item[\code{data}] A `data.frame` or `data.table` containing the input text data.

\item[\code{post\_id}] Character. Column name in `data` containing unique post identifiers. If `NULL`, 
the function attempts to detect it automatically.

\item[\code{time}] Character. Column name in `data` containing timestamps. If `NULL`, the function 
attempts to detect it automatically.

\item[\code{content}] Character. Column name in `data` containing text content to embed. If `NULL`, 
the function attempts to detect it automatically.

\item[\code{model\_name}] The Hugging Face model to use for embeddings (default is the CPU-friendly, multilingual model "sentence-transformers/paraphrase-multilingual-MiniLM-L12-v2").
Other useful, more complex models to consider are, e.g.:
\begin{itemize}

\item{} Social-media specific, multilingual: "Twitter/twhin-bert-base"
\item{} Multilingual, lightweight model, trained on wikipedia: "distilbert/distilbert-base-multilingual-cased"

\end{itemize}


\item[\code{batch\_size}] Integer. Batch size for processing texts (default is 32L).

\item[\code{max\_length}] Integer. Maximum sequence length for tokenization (default is 512L).

\item[\code{use\_fp16}] Logical, whether to use fp16 precision on GPU if available (default is FALSE).

\item[\code{trust\_remote\_code}] Logical, whether to trust remote code from huggingface or not (default is TRUE). Necessary for Alibaba gte.

\item[\code{verbose}] Logical. If `TRUE`, prints progress messages.

\item[\code{uv\_cache\_dir}] Character (optional). Directory used by uv to install python libraries, passed on to initialize\_coorsim.

\item[\code{models\_dir}] Character (optional).  Directory used to cache huggingface models, passed on to initialize\_coorsim.
\end{ldescription}
\end{Arguments}
%
\begin{Value}
A matrix with rownames corresponding to the 'id' column and computed embeddings as columns.
\end{Value}
\HeaderA{initialize\_coorsim}{Initialize the coorsim Python environment}{initialize.Rul.coorsim}
%
\begin{Description}
Declares all required Python packages via \code{\LinkA{py\_require}{py.Rul.require}},
which provisions them using uv into an ephemeral virtual environment.
PyTorch is installed with \code{UV\_TORCH\_BACKEND=auto}, which lets uv
auto-detect NVIDIA (CUDA), AMD (ROCm), and Intel GPUs and pick the correct
wheels — no manual driver querying needed. Pass \code{gpu = FALSE} to force
CPU-only onnxruntime.

The function is called automatically at the start of each coorsim function
and is a no-op after the first successful call within a session.
\end{Description}
%
\begin{Usage}
\begin{verbatim}
initialize_coorsim(gpu = check_gpu(), uv_cache_dir = NULL, models_dir = NULL)
\end{verbatim}
\end{Usage}
%
\begin{Arguments}
\begin{ldescription}
\item[\code{gpu}] Logical. Whether to install \code{onnxruntime-gpu} (CUDA only).
Defaults to auto-detection via \code{check\_gpu()}.

\item[\code{uv\_cache\_dir}] Character (optional). Directory used by uv to install python libraries.

\item[\code{models\_dir}] Character (optional).  Directory used to cache huggingface models.
\end{ldescription}
\end{Arguments}
%
\begin{Details}
Non-torch packages are declared with \code{py\_require()} so that reticulate /
uv can resolve them together. PyTorch is handled by setting
\code{UV\_TORCH\_BACKEND=auto}, which instructs uv to query for CUDA, ROCm,
and Intel GPU drivers and select the matching PyTorch index automatically.

\code{onnxruntime-gpu} only supports CUDA, so the \code{gpu} parameter
reflects NVIDIA presence only; AMD and Intel users get CPU onnxruntime while
still getting GPU-accelerated PyTorch.
\end{Details}
%
\begin{Value}
Invisibly returns \code{TRUE}.
\end{Value}
\HeaderA{label\_communities}{Label Communities Using a Local LLM}{label.Rul.communities}
%
\begin{Description}
This function uses a locally hosted large language model (LLM) via the \pkg{rollama} interface
to generate structured labels and descriptive summaries for user communities. It first labels
each community (or, if large, its individual slices) based on user-level annotations and text,
then optionally aggregates slice-level results in a second LLM pass to produce one coherent
community-level label and description.
\end{Description}
%
\begin{Usage}
\begin{verbatim}
label_communities(
  groups_data,
  max_n_per_slice = 10,
  min_chars = 1,
  max_chars = Inf,
  min_share = NULL,
  model = "gemma3:12b",
  retries = 3,
  retry_trunc = 4000,
  seed = 42,
  temp = 0,
  prompt = NULL,
  system = NULL,
  system_slices = NULL,
  example = NULL,
  example_slices = NULL,
  answer = NULL,
  answer_slices = NULL,
  schema = NULL,
  keep_description_json = FALSE,
  verbose = TRUE
)
\end{verbatim}
\end{Usage}
%
\begin{Arguments}
\begin{ldescription}
\item[\code{groups\_data}] A list containing at least a \code{user\_labels} element,
typically produced by [sample\_user\_text()]. This element must include community membership
and sampled post text for each user.

\item[\code{max\_n\_per\_slice}] Integer. Maximum number of users per slice when splitting large
communities. Default is \code{10}.

\item[\code{min\_chars}] Integer. Minimum character length required for a valid slice input.
Default is \code{1}.

\item[\code{max\_chars}] Integer. Maximum allowed character length per slice (longer input will be truncated).
Default is \code{Inf}.

\item[\code{min\_share}] Optional numeric. Minimum cumulative posting share (\code{user\_share\_comm})
per slice. Default is \code{NULL} (no filter).

\item[\code{model}] Character string. The local LLM model to query via Ollama
(e.g., \code{"gemma3:12b"}). Default is the task-evaluated \code{"gemma3:12b"}.

\item[\code{retries}] Integer. Maximum number of re-query attempts for invalid or unparseable
JSON responses. Default is \code{3}.

\item[\code{retry\_trunc}] Integer. Character limit for truncating text inputs in retry rounds.
Default is \code{4000}.

\item[\code{seed}] Integer. Random seed for reproducibility. Default is \code{42}.

\item[\code{temp}] Numeric. Sampling temperature for model generation
(\code{0} = deterministic). Default is \code{0.0}.

\item[\code{prompt}] Character. User-level prompt text used to instruct the model when labeling
communities (overrides internal \code{prompts\$prompt\_comm}). Optional.

\item[\code{system}] Character. System-level instruction guiding the model's behavior
(overrides internal \code{prompts\$system\_comm}). Optional.

\item[\code{system\_slices}] Character. Optional override for \code{prompts\$system\_comm\_slices},
used for the aggregation step that combines slice-level labels into one community summary.

\item[\code{example}] Character vector containing few-shot example inputs for community-level labeling.
Defaults to \code{examples\$example\_comm\_text}.

\item[\code{example\_slices}] Character vector containing few-shot example inputs for aggregated
slice-level labeling. Defaults to \code{examples\$example\_comm\_slices\_text}.

\item[\code{answer}] Character vector of corresponding example model outputs for community-level labeling.
Defaults to \code{examples\$example\_comm\_answer}.

\item[\code{answer\_slices}] Character vector of example outputs for aggregated slice-level labeling.
Defaults to \code{examples\$example\_comm\_slices\_answer}.

\item[\code{schema}] Optional JSON-like list defining the expected structured output format.
Defaults to \code{schemata\$schema\_comm}.

\item[\code{keep\_description\_json}] Logical. If \code{FALSE} (default), extracts and concatenates
sentences from structured description objects into clean text. If \code{TRUE}, preserves JSON
structure as text. Only relevant when description is returned as a JSON object with
sentence\_1 through sentence\_5 fields.

\item[\code{verbose}] Logical. Whether to print progress messages, retry information,
and aggregation updates. Default is \code{TRUE}.
\end{ldescription}
\end{Arguments}
%
\begin{Details}
If malformed or incomplete JSON responses are returned, the function retries the query up to
a user-specified number of attempts. It supports both single-slice and multi-slice processing,
automatically detecting and aggregating communities that have been split during preprocessing.


The function executes in two main stages:
\begin{enumerate}

\item{} \strong{Slice-level labeling:} If a community exceeds \code{max\_n\_per\_slice} users, it is divided
into smaller slices using [slice\_community\_text()]. Each slice is labeled independently by the LLM.
\item{} \strong{Community-level aggregation:} For multi-slice communities, the slice-level labels and
descriptions are combined into a JSON array and passed to the model again with
\code{system\_comm\_slices} to generate a single coherent label and summary.

\end{enumerate}


JSON responses are validated using \pkg{jsonlite}; malformed responses trigger retries with truncated text
up to \code{retries} times. The function expects a locally running Ollama instance with the specified
model available. Verify setup using [rollama::ping\_ollama()].
\end{Details}
%
\begin{Value}
The same \code{groups\_data} list, updated with a new \code{community\_labels} element.
This is a \pkg{data.table} with one row per labeled community containing:
\begin{description}

\item[\code{community}] Community identifier.
\item[\code{label}] LLM-generated label for the community.
\item[\code{description}] Concise LLM-generated summary of the community's tone, style, and main topics.
\item[\code{lang}] Predominant language inferred by the model.
\item[\code{topic\_1–topic\_5}] Up to five thematic categories assigned by the model.
\item[\code{named\_entity\_1–named\_entity\_5}, \code{sentiment\_1–sentiment\_5}] Named entities and associated sentiments.
\item[\code{pattern\_1–pattern\_3}] Recurring emojis, slogans, or stylistic markers.
\item[\code{incivility}] Whether the majority of users use uncivil or offensive language ("yes"/"no").
\item[\code{elaborate}] Whether the majority of users post elaborate content ("yes"/"no").
\item[\code{confidence}] Model-estimated confidence score between 0 and 1.
\item[\code{text}] Input text provided to the model (either slice or aggregated JSON).
\item[\code{is\_valid\_json}] Logical flag indicating successful JSON validation.
\item[\code{model}, \code{model\_queried\_at}, \code{model\_total\_duration}] Metadata about the model run.

\end{description}

\end{Value}
%
\begin{SeeAlso}
[sample\_user\_text()], [slice\_community\_text()], [rollama::query()], [rollama::make\_query()], [jsonlite::fromJSON()]
\end{SeeAlso}
\HeaderA{label\_users}{Label Users Using a Local LLM}{label.Rul.users}
%
\begin{Description}
This function uses a locally hosted large language model (LLM) through the 
\pkg{rollama} interface to automatically generate structured annotations for users 
based on their sampled posts. It applies a few-shot prompting setup with examples 
and a predefined schema to produce machine-readable JSON output describing each user.
\end{Description}
%
\begin{Usage}
\begin{verbatim}
label_users(
  groups_data,
  model = "llama3.1:8b",
  prompt = NULL,
  system = NULL,
  example = NULL,
  answer = NULL,
  schema = NULL,
  retries = 3,
  retry_trunc = 4000,
  seed = 42,
  temp = 0,
  keep_description_json = FALSE,
  cache_path = NULL,
  verbose = TRUE
)
\end{verbatim}
\end{Usage}
%
\begin{Arguments}
\begin{ldescription}
\item[\code{groups\_data}] A list containing at least the element \code{user\_labels}, 
typically produced by [sample\_user\_text()]. This element must include a column \code{text} 
with sampled user content and an \code{account\_id} column.

\item[\code{model}] Character string. Name of the local LLM model to be queried via Ollama. 
Default is \code{"llama3.1:8b"}.

\item[\code{prompt}] Optional character string providing user-facing task instructions 
(defaults to the internal \code{prompts\$prompt\_user}).

\item[\code{system}] Optional character string defining the system-level instruction to guide model behavior 
(defaults to the internal \code{prompts\$system\_user}).

\item[\code{example}] Optional character vector containing example input texts 
(defaults to \code{examples\$example\_user\_text}).

\item[\code{answer}] Optional character vector containing corresponding example model outputs 
(defaults to \code{examples\$example\_user\_answer}).

\item[\code{schema}] Optional JSON-like list defining the expected output structure 
(defaults to \code{schemata\$schema\_user}).

\item[\code{retries}] Integer. Maximum number of re-query attempts for responses that fail JSON validation. 
Default is \code{3}.

\item[\code{retry\_trunc}] Integer. Character limit for truncating user text during retry attempts to prevent 
context overflow. Default is \code{4000}.

\item[\code{seed}] Integer. Random seed used for model sampling; incremented by 1 on each retry. Default is \code{42}.

\item[\code{temp}] Numeric. Sampling temperature for model generation (0 = deterministic). Default is \code{0.0}.

\item[\code{keep\_description\_json}] Logical. If \code{FALSE} (default), extracts and concatenates sentences from 
structured description objects into clean text. If \code{TRUE}, preserves JSON structure as text. Only relevant 
when description is returned as a JSON object with sentence\_1 through sentence\_5 fields.

\item[\code{cache\_path}] Optional character string; passed on to `rollama::query()` to save responses - to resume after interruption.

\item[\code{verbose}] Logical. Whether to display progress messages and retry information. Default is \code{TRUE}.
\end{ldescription}
\end{Arguments}
%
\begin{Details}
The model should run via an active \emph{Ollama} Docker container with GPU (CUDA) support.
Use `system("docker ps")` to verify that Ollama is running, and ensure that the 
desired model (e.g., `"llama3.1:8b"`) has been pulled via Ollama before executing this function.


The function constructs few-shot prompts by combining examples and instructions, sends them to the 
locally hosted model using [rollama::query()], and validates the returned JSON using \pkg{jsonlite}.
If a response fails validation, it retries up to \code{retries} times with shorter input text 
(\code{retry\_trunc}) and incremented seed values.

The schema guides the model to output structured JSON containing linguistic, emotional, and stylistic annotations 
for each user. Common fields include `description`, `lang`, `topic`, `named\_entities`, `repetitive\_patterns`, 
`emotion\_valence`, `incivility`, `elaborate`, and `confidence`.

Ollama must be accessible locally (e.g., via Docker). Use [rollama::ping\_ollama()] to confirm connectivity.
\end{Details}
%
\begin{Value}
The input \code{groups\_data} list with the \code{user\_labels} element updated to include:
\begin{description}

\item[\code{description}] Concise summary generated by the model.
\item[\code{lang}] Predicted language code (ISO format).
\item[\code{emotion\_valence}] Overall emotional tone ("positive", "negative", "neutral").
\item[\code{incivility}] Presence of uncivil or offensive language ("yes", "no").
\item[\code{elaborate}] Linguistic elaboration category ("simple", "moderate", or "elaborate").
\item[\code{confidence}] Model-assessed confidence (0–1 scale).
\item[\code{topic\_1}–\code{topic\_5}] Detected main topics.
\item[\code{pattern\_1}–\code{pattern\_3}] Repeated emojis, slogans, or stylistic markers.
\item[\code{named\_entity\_1}–\code{named\_entity\_3}, \code{sentiment\_1}–\code{sentiment\_3}] Named entities 
and their corresponding sentiment labels.
\item[\code{is\_valid\_json}] Logical flag indicating valid JSON structure in model output.
\item[\code{response}] Raw model output (JSON string).
\item[\code{model}, \code{model\_queried\_at}, \code{model\_total\_duration}] Metadata for reproducibility.

\end{description}

\end{Value}
%
\begin{SeeAlso}
[sample\_user\_text()], [rollama::query()], [rollama::make\_query()], 
[jsonlite::fromJSON()], [data.table::as.data.table()]
\end{SeeAlso}
\HeaderA{load\_h5\_embeddings}{Load Embeddings from an HDF5 (.h5) File}{load.Rul.h5.Rul.embeddings}
%
\begin{Description}
This function loads precomputed embeddings stored in an `.h5` file and retrieves only the relevant embeddings matching post IDs present in a dataset.
\end{Description}
%
\begin{Usage}
\begin{verbatim}
load_h5_embeddings(path, ids_subset = NULL, verbose = TRUE)
\end{verbatim}
\end{Usage}
%
\begin{Arguments}
\begin{ldescription}
\item[\code{path}] Character. Path to the `.h5` file containing embeddings.

\item[\code{ids\_subset}] Character. Vector of IDs to subset the embedding matrix. Default is `NULL`.

\item[\code{verbose}] Logical. If `TRUE`, displays progress messages. Default is `TRUE`.
\end{ldescription}
\end{Arguments}
%
\begin{Details}
The function checks whether the `.h5` file contains a `metadata/post\_id` dataset and loads the corresponding embeddings for matching post IDs.
\end{Details}
%
\begin{Value}
A numeric matrix where rows correspond to post IDs and columns represent embedding dimensions. If no matching post IDs are found, returns `NULL`.
\end{Value}
\HeaderA{plot\_communities}{Plot Coordinated Communities in a Network}{plot.Rul.communities}
%
\begin{Description}
Visualizes communities of coordinated accounts within a network graph, either as a single overview plot
or as a grid of community-level plots with metadata side panels. Nodes represent accounts, edges reflect
similarity or interaction, and layout and annotations emphasize coordination patterns.
\end{Description}
%
\begin{Usage}
\begin{verbatim}
plot_communities(
  network_data,
  palette_option = "A",
  end_color = 1,
  start_color = 0,
  title_prefix = NULL,
  n_top_nodes = 10,
  label_fontsize = NULL,
  by_community = TRUE,
  ncol = 2,
  expand_hairballs = NULL,
  use_palette = NULL,
  use_order = NULL,
  graph_layout = "stress",
  check_overlap = FALSE,
  add_node_names = TRUE
)
\end{verbatim}
\end{Usage}
%
\begin{Arguments}
\begin{ldescription}
\item[\code{network\_data}] A named list containing at least a simplified graph object (`network\_data\$graph`, as an `igraph`) 
and a node table (`network\_data\$node\_list`, coercible to a tibble). Optionally includes:
- `params`: named list of parameters (e.g., cluster method, thresholds)
- `filter`: optional filter metadata
- `community\_labels`: a data frame with `community`, `label`, `description`, `n\_accs`, and `n\_posts`

\item[\code{palette\_option}] Character. Viridis palette option (e.g., "A", "B", "C", "D").

\item[\code{end\_color}] Numeric. End value for viridis color scale (between 0 and 1).

\item[\code{start\_color}] Numeric. Start value for viridis color scale (between 0 and 1).

\item[\code{title\_prefix}] Optional prefix for the plot title.

\item[\code{n\_top\_nodes}] Integer. Number of top-degree nodes per community to label.

\item[\code{label\_fontsize}] Integer. Base font size for text panels. If NULL, size is auto-scaled.

\item[\code{by\_community}] Logical. If TRUE, create individual plots per community with metadata side panels.

\item[\code{ncol}] Integer. Number of columns in the grid layout when by\_community = TRUE.

\item[\code{expand\_hairballs}] Numeric. Factor to visually expand dense communities.

\item[\code{use\_palette}] Optional palette: function(n) or character vector of hex colors.

\item[\code{use\_order}] Optional vector specifying the order in which communities are plotted.

\item[\code{graph\_layout}] Character. The graph layout algorithm passed
to ggraph::create\_layout(). Defaults to `"stress"`. Other options include "fr", "kk",
"graphopt", and "lgl".

\item[\code{check\_overlap}] Logical. Check overlap when plotting node names? Passed on to `ggraph::geom\_node\_text()`.

\item[\code{add\_node\_names}] Logical. Whether to add node name labels to the plot. Default is `TRUE`.
\end{ldescription}
\end{Arguments}
%
\begin{Value}
A ggplot or patchwork object.
\end{Value}
\HeaderA{plot\_posts}{Plot Coordinated Posts Over Time}{plot.Rul.posts}
%
\begin{Description}
Creates line plots showing coordinated posts over time. The function can
either plot all communities together (colored lines) or produce individual
community-level time series arranged in a grid.
\end{Description}
%
\begin{Usage}
\begin{verbatim}
plot_posts(
  network_data,
  by_community = TRUE,
  unit = "6 hours",
  palette_option = "A",
  start_color = 0,
  end_color = 1,
  use_palette = NULL,
  use_order = NULL,
  title_prefix = NULL,
  ncol = 2,
  include_examples = FALSE,
  n_examples = 2,
  trunc = 100,
  label_fontsize = NULL,
  legend = TRUE,
  seed = 42,
  verbose = TRUE
)
\end{verbatim}
\end{Usage}
%
\begin{Arguments}
\begin{ldescription}
\item[\code{network\_data}] A list containing:
- post\_data: data.table or data.frame with columns account\_id, community, time
- node\_list: community membership table (account\_id, community)
- params, filter: optional lists for caption metadata
- community\_labels: optional data.frame with columns community, label

\item[\code{by\_community}] Logical. If TRUE, plot one panel per community; if FALSE
plot all communities together in one plot.

\item[\code{unit}] Character time unit passed to lubridate::floor\_date().

\item[\code{palette\_option}] Viridis palette option ("A","B","C","D").

\item[\code{start\_color}, \code{end\_color}] Numeric range for viridis palette.

\item[\code{use\_palette}] Optional custom palette function(n) or a character vector
of hex colors (length >= number of communities).

\item[\code{use\_order}] Optional vector specifying the order in which communities
are plotted (and mapped to colors). Values are matched to community ids.

\item[\code{title\_prefix}] Optional prefix for plot title.

\item[\code{ncol}] Number of columns for per-community layout.

\item[\code{include\_examples}, \code{n\_examples}, \code{trunc}] Reserved for future use.

\item[\code{label\_fontsize}] Base font size.

\item[\code{legend}] Logical. Show legend when by\_community = FALSE.

\item[\code{seed}] Random seed (not used yet).

\item[\code{verbose}] Print progress messages.
\end{ldescription}
\end{Arguments}
%
\begin{Value}
A ggplot or patchwork object.
\end{Value}
\HeaderA{prompt\_data}{Prompt and example data for LLM community annotation}{prompt.Rul.data}
\aliasA{examples}{prompt\_data}{examples}
\aliasA{prompts}{prompt\_data}{prompts}
\aliasA{schemata}{prompt\_data}{schemata}
\keyword{datasets}{prompt\_data}
%
\begin{Description}
These objects bundle prompt templates and few-shot examples used for guiding large language models (LLMs) in the `coorsim` package.
All examples are anonymized and rewritten using GPT-4o to standardize tone and protect privacy.
\end{Description}
%
\begin{Usage}
\begin{verbatim}
prompts

examples

schemata
\end{verbatim}
\end{Usage}
%
\begin{Format}
Lists of character vectors.

An object of class \code{list} of length 6.

An object of class \code{list} of length 2.
\end{Format}
%
\begin{Section}{Prompts (`prompts`)}

\begin{description}

\item[prompt\_user] Instruction for generating a user description and language tag in JSON.
\item[prompt\_comm] Instruction for generating a community label and description in JSON.
\item[system\_user] System message instructing how to interpret individual user text.
\item[system\_comm] System message for generating a label and summary from community-level descriptions.
\item[system\_comm\_slices] System message for summarizing community slices by share.

\end{description}

\end{Section}
%
\begin{Section}{Examples (`examples`)}

\begin{description}

\item[example\_user\_text] Simulated user posts in English and German.
\item[example\_user\_answer] Expected output: user description and language.
\item[example\_comm\_text] Concatenated user descriptions per community.
\item[example\_comm\_answer] Expected output: community label and description.
\item[example\_comm\_slices\_text] Pre-labeled community slices with share info.
\item[example\_comm\_slices\_answer] Expected output for slice aggregation.

\end{description}

\end{Section}
%
\begin{Section}{Schemata (`schemata`)}

\begin{description}

\item[schema\_user] JSON-schema for LLM-output to describe users.
\item[schema\_comm] JSON-schema for LLM-output to describe communities.

\end{description}

\end{Section}
\HeaderA{query\_and\_compute\_similarities}{Query Embeddings and Compute Cosine Similarities}{query.Rul.and.Rul.compute.Rul.similarities}
%
\begin{Description}
This function takes a numeric matrix containing embeddings and a list of post ID lists. 
It performs two tasks sequentially to conserve memory:
1. Queries the embeddings for each post ID list.
2. Computes the cosine similarities between the first embedding and all other embeddings within each subset.
\end{Description}
%
\begin{Usage}
\begin{verbatim}
query_and_compute_similarities(m, post_id_lists, threshold)
\end{verbatim}
\end{Usage}
%
\begin{Arguments}
\begin{ldescription}
\item[\code{m}] A NumericMatrix containing the embedding data with row names representing post IDs.

\item[\code{post\_id\_lists}] A List of CharacterVector objects, where each CharacterVector represents a set of post IDs to query in the embedding matrix.

\item[\code{threshold}] A double value representing the threshold for cosine similarity. Only similarities above this threshold are included in the output.
\end{ldescription}
\end{Arguments}
%
\begin{Details}
The function returns a list of data frames where each data frame contains post ID pairs and their cosine similarities that exceed the given threshold.
\end{Details}
%
\begin{Value}
A List of DataFrames. Each data frame contains three columns: "post\_id" (the first post ID in the pair), "post\_id\_y" (the second post ID in the pair), and "similarity" (the cosine similarity between them).
\end{Value}
\HeaderA{query\_and\_compute\_similarities\_tbb}{Query embeddings and compute cosine similarities (group-only, cached)}{query.Rul.and.Rul.compute.Rul.similarities.Rul.tbb}
%
\begin{Description}
Query embeddings and compute cosine similarities (group-only, cached)
\end{Description}
%
\begin{Usage}
\begin{verbatim}
query_and_compute_similarities_tbb(m, post_id_lists, threshold)
\end{verbatim}
\end{Usage}
%
\begin{Arguments}
\begin{ldescription}
\item[\code{m}] NumericMatrix with embeddings; rownames are post IDs.

\item[\code{post\_id\_lists}] List of CharacterVector; each vector is one group.

\item[\code{threshold}] double; only similarities strictly greater than this are returned.
\end{ldescription}
\end{Arguments}
%
\begin{Value}
List of DataFrame (one per input group) with columns: post\_id, post\_id\_y, similarity.
\end{Value}
\HeaderA{replace\_emoji\_with\_name}{Replace Emojis with Descriptive Names}{replace.Rul.emoji.Rul.with.Rul.name}
%
\begin{Description}
Replaces emoji characters in a character vector with their corresponding
descriptive names in `:colon\_syntax:` (e.g., "🔥" becomes ":fire:").
Only emojis present in the input are processed to maximize performance.
\end{Description}
%
\begin{Usage}
\begin{verbatim}
replace_emoji_with_name(text_vec)
\end{verbatim}
\end{Usage}
%
\begin{Arguments}
\begin{ldescription}
\item[\code{text\_vec}] A character vector containing text with emojis to be replaced.
\end{ldescription}
\end{Arguments}
%
\begin{Details}
This function uses the `emoji::emoji\_name` dataset for emoji-to-name mapping
and `stringi` for high-performance detection and substitution.
Only emojis that are actually present in the input are processed, making
the function suitable for large text corpora.
\end{Details}
%
\begin{Value}
A character vector of the same length as `text\_vec`, where all detected
emojis are replaced with their descriptive names in `:colon:` format.
\end{Value}
\HeaderA{sample\_user\_text}{Sample and format user posts for labeling (JSON string)}{sample.Rul.user.Rul.text}
%
\begin{Description}
Samples posts per user and returns a compact JSON-style character string per user.
Character fields are cleaned and truncated to \code{max\_chars}. Empty fields are omitted.
\end{Description}
%
\begin{Usage}
\begin{verbatim}
sample_user_text(
  groups_data,
  user_vars = NULL,
  user_vars_rename = NULL,
  post_vars = NULL,
  post_vars_rename = NULL,
  min_chars = 1,
  max_chars = 1000,
  clean_name = TRUE,
  min_n_posts = 1,
  max_n_posts = 30,
  sampling_ratio_posts = 0.5,
  min_n_users = 1,
  max_n_users = 30,
  sampling_ratio_users = 0.5,
  strat_sample = TRUE,
  seed = 42
)
\end{verbatim}
\end{Usage}
%
\begin{Arguments}
\begin{ldescription}
\item[\code{groups\_data}] Named list with \code{post\_data} and \code{node\_list} (both \code{data.table}).
\code{post\_data} must contain: \code{post\_id}, \code{account\_id}, \code{account\_name}, \code{content}, \code{community}.

\item[\code{user\_vars}] Character vector of user-level columns to include (besides user name). Default \code{NULL}.

\item[\code{user\_vars\_rename}] Named character vector mapping JSON field names to existing user columns.
If \code{NULL}, names in \code{user\_vars} are used as JSON names. Default \code{NULL}.

\item[\code{post\_vars}] Character vector of post-level columns to include per post (besides content). Default \code{NULL}.

\item[\code{post\_vars\_rename}] Named character vector mapping JSON field names to existing post columns.
If \code{NULL}, names in \code{post\_vars} are used as JSON names. Default \code{NULL}.

\item[\code{min\_chars}] Minimum characters required for a cleaned post content to be eligible. Default \code{1}.

\item[\code{max\_chars}] Maximum characters for truncating all cleaned character fields. Use \code{NULL} for no truncation. Default \code{1000}.

\item[\code{clean\_name}] If \code{TRUE}, clean and truncate \code{account\_name} before use as user name. Default \code{TRUE}.

\item[\code{min\_n\_posts}] Minimum number of posts to sample per user. Default \code{1}.

\item[\code{max\_n\_posts}] Maximum number of posts to sample per user. Default \code{30}.

\item[\code{sampling\_ratio\_posts}] Proportion of available posts to sample per user before bounding by min/max. Default \code{0.5}.

\item[\code{min\_n\_users}] Minimum number of users to sample per community; communities below this are dropped. Default \code{1}.

\item[\code{max\_n\_users}] Maximum number of users to sample per community. Default \code{30}.

\item[\code{sampling\_ratio\_users}] Proportion of eligible users to sample per community before bounding by min/max. Default \code{0.5}.

\item[\code{strat\_sample}] If \code{TRUE}, sample users within community with probability proportional to post count; if \code{FALSE}, pick most active. Default \code{TRUE}.

\item[\code{seed}] Integer seed for reproducibility. Default \code{42}.
\end{ldescription}
\end{Arguments}
%
\begin{Details}
All character fields in \code{user\_vars} and \code{post\_vars} are cleaned and truncated
before inclusion. Fields that are \code{NA} or empty after cleaning are omitted from the JSON string.
\end{Details}
%
\begin{Value}
The input \code{groups\_data} with an added \code{user\_labels} \code{data.table}
(one row per sampled \code{account\_id}) with columns:
\begin{itemize}

\item{} \code{account\_id}, \code{community}
\item{} \code{account\_name\_clean} (used as user name)
\item{} \code{text} (compact JSON-style character string with user fields and a list of posts)
\item{} \code{sampled\_post\_ids} (list column)
\item{} \code{approx\_tokens} (approximate token count)
\item{} \code{user\_share\_comm} (user's posting share in community)

\end{itemize}

\end{Value}
\HeaderA{save\_embeddings}{Save Text Embeddings to an HDF5 File}{save.Rul.embeddings}
%
\begin{Description}
This function retrieves text embeddings using a transformer model, processes the 
input data, and stores the resulting embeddings in an HDF5 file along with metadata for efficient 
retrieval and analysis.
\end{Description}
%
\begin{Usage}
\begin{verbatim}
save_embeddings(
  data,
  post_id = NULL,
  time = NULL,
  content = NULL,
  model_name = "sentence-transformers/paraphrase-multilingual-MiniLM-L12-v2",
  batch_size = 32L,
  max_length = 512L,
  use_fp16 = TRUE,
  trust_remote_code = TRUE,
  chunk_size = 512L,
  gzip_level = NULL,
  h5_fileprefix = "embeddings_",
  save_dir = NULL,
  overwrite = TRUE,
  uv_cache_dir = NULL,
  models_dir = NULL,
  verbose = TRUE
)
\end{verbatim}
\end{Usage}
%
\begin{Arguments}
\begin{ldescription}
\item[\code{data}] A `data.frame` or `data.table` containing the input text data.

\item[\code{post\_id}] Character. Column name in `data` containing unique post identifiers. If `NULL`, 
the function attempts to detect it automatically.

\item[\code{time}] Character. Column name in `data` containing timestamps. If `NULL`, the function 
attempts to detect it automatically.

\item[\code{content}] Character. Column name in `data` containing text content to embed. If `NULL`, 
the function attempts to detect it automatically.

\item[\code{model\_name}] Character. Name of the transformer model to use for generating embeddings. 
Defaults to `"sentence-transformers/paraphrase-multilingual-MiniLM-L12-v2"`.

\item[\code{batch\_size}] Integer. Number of texts processed per batch when generating embeddings. 
Defaults to `32L`.

\item[\code{max\_length}] Integer. Maximum token length for transformer-based embedding generation. 
Defaults to `512L`.

\item[\code{use\_fp16}] Logical. Whether to use half-precision floating point (`fp16`) for faster 
computation on GPUs. Defaults to `TRUE`.

\item[\code{trust\_remote\_code}] Logical, whether to trust remote code from huggingface or not (default is TRUE). 
Necessary for Alibaba gte.

\item[\code{chunk\_size}] Integer. The number of observations per chunk when storing embeddings in the 
HDF5 file. Defaults to `512L`.

\item[\code{gzip\_level}] Integer. gzip compression level (1-9) used when storing the HDF5 file. 1: Little compression, fast. 9: High compression, slow.
Default is NULL.

\item[\code{h5\_fileprefix}] Character. Prefix for the generated HDF5 filename. Defaults to `"embeddings\_"`.

\item[\code{save\_dir}] Character. Directory where the HDF5 file will be saved. If `NULL`, defaults to the 
working directory.

\item[\code{overwrite}] Logical. Whether to overwrite an existing HDF5 file.

\item[\code{uv\_cache\_dir}] Character (optional). Directory used by uv to install python libraries, passed on to initialize\_coorsim.

\item[\code{models\_dir}] Character (optional).  Directory used to cache huggingface models, passed on to initialize\_coorsim.

\item[\code{verbose}] Logical. If `TRUE`, prints progress messages. Defaults to `TRUE`.
\end{ldescription}
\end{Arguments}
%
\begin{Details}
The function performs the following steps:
\begin{itemize}

\item{} Ensures `data` contains the required columns (`post\_id`, `time`, `content`), renaming them if necessary.
\item{} Removes duplicate entries based on `post\_id`.
\item{} Converts `time` to a UNIX timestamp format for interoperability.
\item{} Sorts the dataset by `time` to enable fast querying.
\item{} Generates text embeddings using the specified transformer model.
\item{} Stores embeddings and metadata (post IDs and timestamps) in an HDF5 file using chunked storage.

\end{itemize}


The HDF5 file structure includes:
\begin{itemize}

\item{} `"metadata/time"`: Vector of timestamps (UNIX format).
\item{} `"metadata/post\_id"`: Vector of post IDs.
\item{} `"embeddings"`: Matrix of text embeddings stored in chunks.

\end{itemize}


The chunked storage format enables efficient retrieval of embeddings by time range.
\end{Details}
%
\begin{Value}
Invisibly returns `NULL`. The function is executed for its side effects of generating 
and storing text embeddings.
\end{Value}
\HeaderA{set\_num\_threads}{Set number of threads}{set.Rul.num.Rul.threads}
%
\begin{Description}
Set number of threads
\end{Description}
%
\begin{Usage}
\begin{verbatim}
set_num_threads(threads)
\end{verbatim}
\end{Usage}
%
\begin{Arguments}
\begin{ldescription}
\item[\code{threads}] Integer, setting the number of threads for parallel processing.
\end{ldescription}
\end{Arguments}
\HeaderA{slice\_community\_text}{Slice Community-Level User Descriptions for LLM Annotation}{slice.Rul.community.Rul.text}
%
\begin{Description}
This function prepares community-level text slices for LLM-based labeling by
aggregating user-level descriptions sampled via [sample\_user\_text()]. It cleans,
truncates, and filters user descriptions, inserts the user's name into the JSON
structure, and aggregates the resulting JSON objects into compact community-level
text slices for annotation.
\end{Description}
%
\begin{Usage}
\begin{verbatim}
slice_community_text(
  groups_data,
  max_n_per_slice = 10,
  min_chars = 20,
  max_chars = 500,
  min_share = NULL,
  drop_fields = c("confidence", "emotion_valence"),
  verbose = TRUE,
  seed = 42
)
\end{verbatim}
\end{Usage}
%
\begin{Arguments}
\begin{ldescription}
\item[\code{groups\_data}] A list that must include the element `user\_labels`, typically
created by [sample\_user\_text()]. This table must include at least the columns
`community`, `user\_share\_comm`, `response`, and `account\_name\_clean`.

\item[\code{max\_n\_per\_slice}] Integer. Maximum number of user descriptions included
per text slice (per community). Default is `10`.

\item[\code{min\_chars}] Integer. Minimum number of characters required in a user
description to be retained. Default is `20`. *(Currently not enforced in filtering,
but reserved for future use.)*

\item[\code{max\_chars}] Integer. Maximum number of characters retained per description
after cleaning. Default is `500`. *(Currently not applied in truncation, but included
for consistency with [sample\_user\_text()].)*

\item[\code{min\_share}] Optional numeric. Minimum cumulative `user\_share\_comm` required
for a slice to be retained. Default is `NULL` (no filtering).

\item[\code{drop\_fields}] Character vector of JSON fields to exclude from each user’s
description (e.g., `"confidence"`, `"emotion\_valence"`). Default is `c("confidence", "emotion\_valence")`.

\item[\code{verbose}] Logical. If `TRUE`, prints the number of returned community slices.
Default is `TRUE`.

\item[\code{seed}] Integer. Random seed for reproducibility. Default is `42`.
\end{ldescription}
\end{Arguments}
%
\begin{Details}
This function aggregates user-level JSON outputs (e.g., from local LLM-based
annotations) into structured, slice-level JSON documents suitable for community
labeling. It automatically embeds user names into the JSON under `"name"`, and
drops unwanted fields defined in `drop\_fields`. Each slice corresponds to the
top `max\_n\_per\_slice` users within a community, ranked by `user\_share\_comm`.
\end{Details}
%
\begin{Value}
A `data.table` with one row per community–slice combination containing:
\begin{itemize}

\item{} \code{community} – Community identifier.
\item{} \code{slice} – Slice index within the community.
\item{} \code{n} – Number of user descriptions included in the slice.
\item{} \code{sum\_user\_share\_comm} – Sum of user shares within the slice.
\item{} \code{text\_slice} – JSON-formatted string of concatenated user descriptions,
including user names and excluding dropped fields.

\end{itemize}

\end{Value}
%
\begin{SeeAlso}
[sample\_user\_text()], [label\_communities()], [jsonlite::toJSON()], [stringr::str\_replace()]
\end{SeeAlso}
\HeaderA{unpack\_description\_json}{Convert JSON Description to Plain Text}{unpack.Rul.description.Rul.json}
%
\begin{Description}
Unpacks a description column stored as JSON text (from keep\_description\_json = TRUE) 
into concatenated plain text by extracting and joining sentence\_1 through sentence\_5.
\end{Description}
%
\begin{Usage}
\begin{verbatim}
unpack_description_json(description_json)
\end{verbatim}
\end{Usage}
%
\begin{Arguments}
\begin{ldescription}
\item[\code{description\_json}] Character vector containing JSON-formatted descriptions
\end{ldescription}
\end{Arguments}
%
\begin{Value}
Character vector with concatenated sentence text
\end{Value}
%
\begin{Examples}
\begin{ExampleCode}
## Not run: 
# Convert a data frame column
df$description <- unpack_description_json(df$description)

# Or use with dplyr
df <- df |> dplyr::mutate(description = unpack_description_json(description))

## End(Not run)

\end{ExampleCode}
\end{Examples}
\HeaderA{z\_standardize}{Z-score Standardization}{z.Rul.standardize}
%
\begin{Description}
This function performs Z-score standardization (also called normalization) on a numeric vector. 
It scales the input data so that it has a mean of 0 and a standard deviation of 1.
\end{Description}
%
\begin{Usage}
\begin{verbatim}
z_standardize(x)
\end{verbatim}
\end{Usage}
%
\begin{Arguments}
\begin{ldescription}
\item[\code{x}] A numeric vector to be standardized. Missing values (NA) will be ignored in the calculation.
\end{ldescription}
\end{Arguments}
%
\begin{Details}
The function calculates the mean and standard deviation of the input vector, ignoring any 
missing values (`NA`). It then subtracts the mean from each value and divides the result by the standard deviation.
\end{Details}
%
\begin{Value}
A numeric vector of the same length as `x`, where each value is standardized to have a mean of 0 
and a standard deviation of 1. Any missing values (NA) in `x` are preserved.
\end{Value}
%
\begin{Examples}
\begin{ExampleCode}
# Example usage
x <- c(1, 2, 3, 4, 5)
z_standardize(x)

\end{ExampleCode}
\end{Examples}
\printindex{}
\end{document}
